\normalsize\justifying
\setlength{\parskip}{2mm}
\setcounter{page}{1}
\section*{\Huge Introduction}
\addcontentsline{toc}{chapter}{Introduction}

``Quantum mechanics is the description of the behavior of matter and light in all details and, in particular, of the happenings on an atomic scale'' said \textit{R.Feynman} in one of his lectures~\cite{Feynman}. In this framework, position is no longer a single well-defined value prior to measurement, instead, position is described by a probability density all across space. Consequently, the notion of a trajectory as a continuous succession of points in spacetime is lost in quantum dynamics. The definition of trajectory will therefore be addressed, but one question remains paramount when mentioning systems: how can we even define a quantum system properly?\par

In the ideal case, systems are treated as \underline{closed}, i.e., isolated from any external environment and interacting with nothing at all. This modelling allows quantum mechanics to describe the dynamics of systems with unitary and deterministic terms (the state at time $t$ is determined solely by knowing the state at time $t_0$) using the Schrödinger equation.~\cite{CT}However, in the physical world, systems do ``feel'' their surroundings, \textit{i.e.}, they interact with their environment. This interaction should be taken into account to obtain a more complete and physically relevant representation. In doing so, we consider the so-called \textit{Open Quantum System} (OQS) theory~\cite{Breuer}.\par 

Other useful theories have been developed through the last decades, allowing quantum mechanics' description to be more and more complete. One of them is the Quantum Measurement Theory (QMT)~\cite{Jacobs}, where the concept of measurement is fully described. Measuring a quantum system has a tremendous impact on its properties. Furthermore, since any experimental probe inevitably interacts with the system, it is crucial to figure out the nature of this interaction. Measuring a system without impacting its properties leads to the concept of \underline{weak measurements}, that we will develop further below.\par

By studying QMT, one can thus notice that measuring a quantum system alters it. The measurement usually interacts so strongly with the system that it collapses into a state --- an eigenstate of the measured observable --- leading to a loss of the system's properties of interest. For example it can lose its entanglement\footnote{
    $\ket{\psi}_{AB}=\frac{\ket{00}+\ket{11}}{\sqrt{2}}\underset{\text{in A}}{\overset{\text{measuring 0}}{\longrightarrow}}\ket{00}_{AB}=\ket{0}_A\otimes\ket{0}_B$ a product state (thus non-entangled).
}. Along the same lines, quantum systems can be in superposition of states. To illustrate this property, we can think about the \textit{Schrödinger's cat} thought experiment, which corresponds formally to a superposition of `the cat is alive' and `the cat is dead' states prior to measurement. Measuring the state of the cat will determine its fate, alive or dead. It is then ``stuck'' in the state measured and cannot be expressed as a superposition again.
The idea of \textbf{weak measurements} is to \textbf{measure the environment} coupled with the system. Instead of altering dramatically the state of the system, we extract pieces of information about it. The measurements can even be conducted continuously, maximizing the amount of data that an observer can obtain. These are the global ideas behind \textbf{Quantum Continuous Measurements} (QCM)~\cite{Albarelli}.

Practically, this work will be organized as follows: the first chapter covers an academical progression from the basics of quantum mechanics up to the aforementioned theories. A second chapter will cover the \underline{photodetection} and \underline{homodyne detection} of the environment of open quantum systems. The final section of this second chapter presents simulation results for both the qubit and the harmonic oscillator, comparing the two detection methods on each of them.